\subsection{并集与交集的补集}

命题5。对于每一个族 \((\mathbf{X}_t)_{t\in I}\) 对于集合E的子集族，我们有

\[
\mathbb {G} _ {\mathrm{E}} \left(\bigcup_ {i \in \mathrm{I}} \mathrm{X} _ {i}\right) = \bigcap_ {i \in \mathrm{I}} \left(\mathbb {G} _ {\mathrm{E}} \mathrm{X} _ {i}\right), \quad \mathbb {G} _ {\mathrm{E}} \left(\bigcap_ {i \in \mathrm{I}} \mathrm{X} _ {i}\right) = \bigcup_ {i \in \mathrm{I}} \left(\mathbb {G} _ {\mathrm{E}} \mathrm{X} _ {i}\right).
\]

让 \(x \in \complement_{\mathbf{E}} \left( \bigcup_{i \in \mathbf{I}} \mathbf{X}_i \right)\) 。然后 \(x \in E\) 并且，对于每一个 \(i \in I\) , \(x \notin \mathbf{X}_i\) ，因此 \(x \in \complement_{\mathbf{E}} \mathbf{X}_i\) ；因此

\[
x \in \bigcap_ {i \in I} \left(\complement_ {E} X _ {i}\right).
\]

反之，设 \(x\) 为 \(\bigcap_{i \in I} (\complement_E X_i)\) ；根据交的定义，我们有 \(x \in E\) 。此外，如果我们有 \(x \in \bigcup_{i \in I} X_i\) ，那么将存在一个指标 \(x \in I\) 使得 \(x \in X_x\) ，这与假设相矛盾 \(x \in \bigcap_{i \in I} (\complement_E X_i)\) ；因此

\[
x \in \mathbb {C} _ {\mathbf {E}} \left(\bigcup_ {i \in \mathbf {I}} \mathrm{X} _ {i}\right).
\]

这证明了第一个公式；第二个公式是直接推论，鉴于关系式 \(\mathbf{G}_{\mathbf{E}}(\mathbf{G}_{\mathbf{E}}\mathbf{X})=\mathbf{X}\) 对 E 的每个子集 X 成立。

